\scanpage{24}

Je ne vais pas faire l'explicitation ici.

2.6. Le rev. quadratique \inserted{de $X$} associé à $\mathcal A_Q$ mérite de s'appeler l'\uline{invariant d'Arf} de la forme quadratique binaire $Q$. On définira un inv. d'Arf pour \uline{\unsure{tte|les} forme} quadratique scalaire sur un module loc. libre de rang donné (mais en distinguant les cas du rang pair et du rang impair, pour la nature de cet invariant \ldots).

3. \uline{Théorème récapitulatif}

Equivalences entre \deleted{\unsure{deux}} catégories fibrées sur \inserted{la catégorie} $\mathrm{Sch}$\deleted{\unsure{ém}} des schémas $X$ :

\textcircled{1} Structures de similitudes binaires \inserted{str. fidèles} sur $X$, i.e. des systèmes

\[
(E,\deleted{\unsure{Q}}\Delta),
\]

où $E$ est un \deleted{\unsure{M}}module loc. libre de rang $2$ sur $X$, et

\[
\Delta\subset\underline{\operatorname{Quad}}(E)
\]

un sous-\deleted{Module \unsure{inver.}}\inserted{fibré vectoriel} \uline{inversible} de $\underline{\operatorname{Quad}}(E)$, i.e. un sous-module inversible \inserted{et} loc. fact. direct (ou encore, tel que $\Delta\longrightarrow\underline{\operatorname{Quad}}\,E$ soit univ. injectif \unsure{--}).
% archive sheet number: 23
